Weather and climate models cannot resolve every process, so unresolved ``subgrid'' effects are replaced by parametrisations. Data-driven parametrisations learn this map from high-resolution simulation: Rasp et al.\ trained a deep network on a multiscale model and obtained stable multi-year prognostic runs that reproduced the mean climate and aspects of variability~\cite{rasp2018deep}. Whether a closure that is accurate offline is also good when coupled to the resolved model is a separate question. Rasp noted that online simulations with offline-trained machine-learning parametrisations in Earth system models were frequently plagued by instabilities and biases, proposed coupled online learning as a remedy, and illustrated the concept in the Lorenz-96 model~\cite{rasp2019coupled}; Frezat et al.\ train subgrid parametrisations online for non-differentiable solvers by means of a neural emulator~\cite{frezat2023gradient}. Stochastic closures are the other classical approach: the two-scale Lorenz-96 system~\cite{lorenz2006predictability} has been used to study the effects of stochastic parametrisations~\cite{wilks2005effects} and their link to model uncertainty~\cite{arnold2013stochastic}, and learned stochastic variants use GANs~\cite{gagne2019machine} or recurrent networks with red noise~\cite{parthipan2022using}.

These studies differ in their system settings, input features and scores, so it is hard to read off from them how the three most basic closures compare. We use the two-scale Lorenz-96 system as a cheap, idealised testbed and ask: with identical training data, inputs and metrics, how do a regression polynomial, a small MLP and a polynomial with AR(1) noise differ in (i) short-term forecast skill, (ii) long-rollout stability and (iii) the climate of the slow variables, and does the ranking change with scale separation and a forcing shift? This is a controlled small-scale comparison, not a new method. Our contributions are: a fully specified protocol with a sampling-noise floor from an independent truth run; five-seed results at two scale-separation regimes with the hypotheses and tests fixed before the runs; and ablations on polynomial degree, noise amplitude and memory, MLP input width and forcing shift, including cases where a closure fails.
